\chapter{Distributions and Sobolev Spaces}\label{distributions}
\chapterstatus{complete}

\section{Introduction}\label{distributions:section-introduction}

To solve $\Delta u = f$ one first finds a weak solution in a space of distributions and then proves that it is smooth. This chapter
sets up the two tools that make this work: the Fourier transform, which turns differentiation into multiplication, and the Sobolev
spaces $H^s$, in which the elliptic estimate (Theorem~\ref{distributions:theorem-elliptic-estimate}) and Rellich's compactness
theorem (Theorem~\ref{distributions:theorem-rellich}) hold. For operators with variable coefficients the estimate and the regularity theorem are proved by freezing the coefficients and by difference
quotients (Theorems~\ref{distributions:theorem-local-estimate} and~\ref{distributions:theorem-local-regularity}), with Peetre's inequality and the behaviour of $H^s$ under multiplication as tools. The Poincar\'e inequalities, which bound a function by its gradient once the constants are
excluded, give the simplest existence theorem, for the Dirichlet problem (Subsection~\ref{distributions:subsection-poincare}). On a compact manifold these local statements are patched together with
partitions of unity to prove the Hodge theorem (Chapter~\ref{hodge-theorem}). References: \cite{hormander-pdo1},
\cite[Chapters 6--7]{rudin-fa}, \cite{brezis-functional}, \cite[Chapter 6]{warner-foundations}.

\noindent\textbf{Prerequisites.} Chapter~\ref{functional-analysis} (Functional Analysis).

Multi-index notation is as in Theorem~\ref{real-analysis:theorem-taylor-several}; we write $D^\alpha = \partial^\alpha$ and $\langle \xi \rangle = (1 + |\xi|^2)^{1/2}$.

\section{Distributions}\label{distributions:section-distributions}

\begin{definition}\label{distributions:definition-distribution}
Let $U \subseteq \RR^n$ be open and $\mathcal{D}(U) = C_c^\infty(U)$. A sequence $\varphi_k \to \varphi$ in $\mathcal{D}(U)$ means: all $\varphi_k$ have support in a fixed compact subset of $U$, and
$D^\alpha\varphi_k \to D^\alpha\varphi$ uniformly for every $\alpha$. A \emph{distribution} is a linear map $T \colon \mathcal{D}(U) \to \CC$ such that $T(\varphi_k) \to T(\varphi)$ whenever
$\varphi_k \to \varphi$; equivalently, for every compact $K \subseteq U$ there are $C$ and $N$ with $|T(\varphi)| \leq C\sum_{|\alpha| \leq N}\sup|D^\alpha\varphi|$ for all $\varphi$ supported in $K$. The space of
distributions is $\mathcal{D}'(U)$, and $T_k \to T$ means $T_k(\varphi) \to T(\varphi)$ for every $\varphi$.
\end{definition}

\begin{example}\label{distributions:example-distributions}
\begin{enumerate}
\item Every $f \in L^1_{\mathrm{loc}}(U)$ defines $T_f(\varphi) = \int f\varphi$, and $T_f = T_g$ if and only if $f = g$ a.e.; so $\mathcal{D}'(U)$ contains the locally integrable functions.
\item The \emph{Dirac distribution} $\delta_a(\varphi) = \varphi(a)$ is not of this form: if $\delta_0 = T_f$, then $\int f\varphi = 0$ for all $\varphi$ supported away from $0$, so $f = 0$ a.e., while
$\delta_0 \neq 0$.
\item $\operatorname{pv}\frac1x(\varphi) = \lim_{\varepsilon \to 0}\int_{|x| > \varepsilon}\frac{\varphi(x)}{x}\,dx$ is a distribution on $\RR$ of order $1$.
\end{enumerate}
\end{example}

\begin{definition}\label{distributions:definition-derivative}
For $T \in \mathcal{D}'(U)$, $\alpha$ a multi-index and $g \in C^\infty(U)$, set $(D^\alpha T)(\varphi) = (-1)^{|\alpha|}T(D^\alpha\varphi)$ and $(gT)(\varphi) = T(g\varphi)$. The \emph{support} of $T$ is the
complement of the largest open set on which $T$ vanishes.
\end{definition}

\begin{proposition}\label{distributions:proposition-derivative}
These are distributions; $D^\alpha D^\beta T = D^{\alpha+\beta}T$; for $f \in C^k(U)$ the distributional derivatives of order $\leq k$ are the classical ones; and differentiation is continuous: $T_k \to T$
implies $D^\alpha T_k \to D^\alpha T$.
\end{proposition}

\begin{proof}
Continuity is clear from the definition, and the formula for $D^\alpha T$ is forced by integration by parts, which for $f \in C^k$ and $\varphi \in C_c^\infty$ gives
$\int (D^\alpha f)\varphi = (-1)^{|\alpha|}\int f D^\alpha\varphi$ (Fubini's theorem and the fundamental theorem of calculus in one variable at a time,
Proposition~\ref{real-analysis:proposition-parameter-integrals}). The other statements are immediate.
\end{proof}

\begin{example}\label{distributions:example-heaviside}
The Heaviside function $H = \mathbf{1}_{[0,\infty)}$ on $\RR$ has $H' = \delta_0$, since $-\int_0^\infty\varphi' = \varphi(0)$. Hence $\delta_0' (\varphi) = -\varphi'(0)$, and the function $\log|x|$ has
distributional derivative $\operatorname{pv}\frac1x$.
\end{example}

\begin{proposition}\label{distributions:proposition-zero-derivative}
Let $I \subseteq \RR$ be an open interval. A distribution $T$ on $I$ with $T' = 0$ is a constant function: $T = T_c$ for some $c \in \CC$.
\end{proposition}

\begin{proof}
Fix $\psi \in \mathcal{D}(I)$ with $\int\psi = 1$ and put $c = T(\psi)$. For $\varphi \in \mathcal{D}(I)$, the function $\varphi - (\int\varphi)\psi$ has integral $0$ and compact support in $I$, so it is the derivative of
$\chi(x) = \int_{-\infty}^x(\varphi - (\int\varphi)\psi)$, which also has compact support in $I$ (it vanishes to the left of both supports, and to the right because the total integral is $0$). Hence
$T(\varphi) - (\int\varphi)c = T(\chi') = -T'(\chi) = 0$, that is, $T(\varphi) = \int c\varphi$.
\end{proof}

\begin{proposition}\label{distributions:proposition-approximate-identity}
Let $\rho \in L^1(\RR^n)$ with $\int\rho = 1$ and $\rho_\varepsilon(x) = \varepsilon^{-n}\rho(x/\varepsilon)$. Then $T_{\rho_\varepsilon} \to \delta_0$ in $\mathcal{D}'(\RR^n)$ as $\varepsilon \to 0$.
\end{proposition}

\begin{proof}
For $\varphi \in \mathcal{D}(\RR^n)$, the substitution $x = \varepsilon y$ gives $\int\rho_\varepsilon\varphi - \varphi(0) = \int\rho(y)(\varphi(\varepsilon y) - \varphi(0))\,dy$. The integrand tends to $0$ pointwise and is bounded by
$2\sup|\varphi|\,|\rho|$, so the integral tends to $0$ by dominated convergence (Theorem~\ref{measure-theory:theorem-dominated-convergence}).
\end{proof}

\section{The Fourier transform}\label{distributions:section-fourier}

\begin{definition}\label{distributions:definition-schwartz}
The \emph{Schwartz space} $\mathcal{S}(\RR^n)$ consists of the $f \in C^\infty(\RR^n)$ with $p_{\alpha\beta}(f) = \sup_x|x^\alpha D^\beta f(x)| < \infty$ for all $\alpha, \beta$; convergence in $\mathcal{S}$ means
convergence in all these seminorms. Its continuous dual $\mathcal{S}'(\RR^n)$ is the space of \emph{tempered distributions}. For $f \in \mathcal{S}$,
\[
\hat f(\xi) = \mathcal{F}f(\xi) = \int_{\RR^n}f(x)e^{-2\pi i x \cdot \xi}\,dx .
\]
\end{definition}

\begin{theorem}\label{distributions:theorem-fourier-schwartz}
$\mathcal{F}$ maps $\mathcal{S}$ continuously to $\mathcal{S}$, with $\mathcal{F}(D^\alpha f)(\xi) = (2\pi i\xi)^\alpha\hat f(\xi)$ and $\mathcal{F}(x^\alpha f)(\xi) = (-2\pi i)^{-|\alpha|}D^\alpha\hat f(\xi)$.
Moreover:
\begin{enumerate}
\item $\mathcal{F}(e^{-\pi|x|^2}) = e^{-\pi|\xi|^2}$;
\item (Inversion) $f(x) = \int\hat f(\xi)e^{2\pi ix\cdot\xi}\,d\xi$, so $\mathcal{F}$ is a bijection of $\mathcal{S}$ with inverse $\mathcal{F}^{-1}g(x) = \hat g(-x)$;
\item (Plancherel) $\int f\bar g = \int\hat f\overline{\hat g}$; in particular $\|\hat f\|_2 = \|f\|_2$, and $\mathcal{F}$ extends to a unitary operator on $L^2(\RR^n)$.
\end{enumerate}
\end{theorem}

\begin{proof}
Differentiation under the integral sign and integration by parts give the two formulas, and they imply that $\mathcal{F}f$ is smooth and rapidly decreasing, with continuity of $\mathcal{F}$ in the seminorms.

(1) For $n = 1$ let $g(\xi) = \int e^{-\pi x^2}e^{-2\pi ix\xi}dx$. Differentiating under the integral and integrating by parts, $g'(\xi) = -2\pi\xi g(\xi)$, so $g(\xi) = g(0)e^{-\pi\xi^2}$, and $g(0) = 1$ by
Example~\ref{measure-theory:example-gaussian}. The general case follows by Fubini's theorem, since $e^{-\pi|x|^2} = \prod_je^{-\pi x_j^2}$.

(2) For $f, g \in \mathcal{S}$, Fubini's theorem gives the multiplication formula $\int\hat fg = \int f\hat g$. Apply it with $g(\xi) = e^{2\pi ix\cdot\xi}e^{-\pi\varepsilon^2|\xi|^2}$, whose Fourier transform is
$\varepsilon^{-n}e^{-\pi|x - y|^2/\varepsilon^2}$ by (1) and the scaling and translation rules. The right-hand side is a convolution of $f$ with an approximate identity, which converges to $f(x)$ as
$\varepsilon \to 0$ by uniform continuity, and the left-hand side converges to $\int\hat f(\xi)e^{2\pi ix\cdot\xi}d\xi$ by dominated convergence.

(3) Let $h = \overline{\hat g}$; then $\hat h = \bar g$ by (2) and $\overline{\mathcal{F}\bar u} = \mathcal{F}^{-1}u$. The multiplication formula with $f$ and $h$ gives $\int\hat f\overline{\hat g} = \int f\bar g$.
The space $\mathcal{S}$ is dense in $L^2$: it contains $C_c^\infty$, and every continuous function with compact support is a uniform limit of its convolutions with a smooth approximate identity
(Exercise~\ref{measure-theory:exercise-convolution}), which lie in $C_c^\infty$, while such functions are dense in $L^2$ by Proposition~\ref{measure-theory:proposition-density}. Hence $\mathcal{F}$ extends
uniquely to an isometry of $L^2$,
which is surjective because its image contains the dense subspace $\mathcal{S}$.
\end{proof}

\begin{proposition}\label{distributions:proposition-fourier-l1}
For $f \in L^1(\RR^n)$ the integral $\hat f(\xi) = \int f(x)e^{-2\pi ix\cdot\xi}\,dx$ defines a continuous function with $\|\hat f\|_\infty \leq \|f\|_1$ and $\hat f(\xi) \to 0$ as $|\xi| \to \infty$ (the
\emph{Riemann--Lebesgue lemma}). For $f, g \in L^1$, $\widehat{f * g} = \hat f\,\hat g$.
\end{proposition}

\begin{proof}
The bound is clear, and continuity follows from dominated convergence. For $f \in \mathcal{S}$, $\hat f \in \mathcal{S}$ tends to $0$ at infinity. In general choose $f_k \in \mathcal{S}$ with $\|f - f_k\|_1 \to 0$
($C_c^\infty$ is dense in $L^1$, by the argument in the proof of Theorem~\ref{distributions:theorem-fourier-schwartz}(3)); then $\hat f_k \to \hat f$ uniformly, and a uniform limit of functions tending to
$0$ at infinity tends to $0$ at infinity. For the convolution, $f * g \in L^1$ (Exercise~\ref{measure-theory:exercise-convolution}), and by Fubini's theorem
$\int\int f(x - y)g(y)e^{-2\pi ix\cdot\xi}\,dy\,dx = \int g(y)e^{-2\pi iy\cdot\xi}\int f(x - y)e^{-2\pi i(x - y)\cdot\xi}\,dx\,dy = \hat f(\xi)\hat g(\xi)$.
\end{proof}

\begin{example}\label{distributions:example-fourier-computations}
\begin{enumerate}
\item The indicator of $[-\frac12, \frac12]$ has $\hat f(\xi) = \int_{-1/2}^{1/2}e^{-2\pi ix\xi}\,dx = \frac{\sin\pi\xi}{\pi\xi}$, which is continuous and tends to $0$, but is not integrable: by the inversion theorem it
cannot be, since $f$ is not continuous.
\item $f(x) = e^{-2\pi|x|}$ on $\RR$: $\hat f(\xi) = 2\int_0^\infty e^{-2\pi x}\cos(2\pi x\xi)\,dx = 2\operatorname{Re}\frac{1}{2\pi(1 - i\xi)} = \frac{1}{\pi(1 + \xi^2)}$. A kink of $f$ at $0$ gives a decay
$|\xi|^{-2}$ only; smoothness of $f$ corresponds to decay of $\hat f$, and decay of $f$ to smoothness of $\hat f$.
\item Scaling and translation: $\widehat{f(\lambda x)}(\xi) = \lambda^{-n}\hat f(\xi/\lambda)$ and $\widehat{f(x - a)}(\xi) = e^{-2\pi ia\cdot\xi}\hat f(\xi)$. With
Theorem~\ref{distributions:theorem-fourier-schwartz}(1), $\widehat{e^{-\pi t|x|^2}} = t^{-n/2}e^{-\pi|\xi|^2/t}$ for $t > 0$: the narrower the Gaussian, the wider its transform.
\end{enumerate}
\end{example}

\begin{proposition}[Poisson summation]\label{distributions:proposition-poisson}
For $f \in \mathcal{S}(\RR)$, $\sum_{n \in \ZZ}f(x + n) = \sum_{k \in \ZZ}\hat f(k)e^{2\pi ikx}$ for all $x$; in particular $\sum_nf(n) = \sum_k\hat f(k)$.
\end{proposition}

\begin{proof}
The periodisation $P(x) = \sum_nf(x + n)$ converges uniformly with all derivatives on compact sets, as $f$ and its derivatives decay faster than $|x|^{-2}$; so $P$ is a smooth function on $\RR/\ZZ$. Its
Fourier coefficients are $\int_0^1P(x)e^{-2\pi ikx}\,dx = \sum_n\int_n^{n+1}f(y)e^{-2\pi iky}\,dy = \hat f(k)$. These decay rapidly, so the Fourier series converges absolutely and uniformly to a continuous
function, which equals $P$ almost everywhere, being its $L^2$ limit (Example~\ref{functional-analysis:example-fourier}), hence everywhere.
\end{proof}

For $f(x) = e^{-\pi tx^2}$ this gives the transformation formula of the theta function, $\sum_ne^{-\pi tn^2} = t^{-1/2}\sum_ke^{-\pi k^2/t}$ for $t > 0$, which relates large and small $t$ and is
the starting point of the functional equation of the Riemann zeta function.

\begin{proposition}[Heisenberg's uncertainty principle]\label{distributions:proposition-heisenberg}
For $f \in \mathcal{S}(\RR)$, $\|f\|_2^2 \leq 4\pi\|xf\|_2\,\|\xi\hat f\|_2$. So $f$ and $\hat f$ cannot both be concentrated near $0$.
\end{proposition}

\begin{proof}
Integrating by parts, $\|f\|_2^2 = \int|f|^2\,dx = -\int x\frac{d}{dx}|f|^2\,dx = -2\operatorname{Re}\int x\bar ff^{\prime}\,dx \leq 2\|xf\|_2\|f^{\prime}\|_2$ by Cauchy--Schwarz, and
$\|f^{\prime}\|_2 = \|2\pi\xi\hat f\|_2$ by Plancherel and Theorem~\ref{distributions:theorem-fourier-schwartz}.
\end{proof}

Equality holds for the Gaussians $e^{-\pi tx^2}$; after normalising $\|f\|_2 = 1$, the inequality says that the variances of $|f|^2$ and $|\hat f|^2$ have product at least $1/(16\pi^2)$, the mathematical
content of the uncertainty relation of quantum mechanics.

\begin{definition}\label{distributions:definition-fourier-tempered}
For $T \in \mathcal{S}'$, $\hat T(\varphi) = T(\hat\varphi)$. By Theorem~\ref{distributions:theorem-fourier-schwartz} this defines a linear bijection of $\mathcal{S}'$ extending the transform on $\mathcal{S}$
and on $L^2$, with $\mathcal{F}(D^\alpha T) = (2\pi i\xi)^\alpha\hat T$.
\end{definition}

\begin{example}\label{distributions:example-fourier-delta}
$\hat\delta_0 = 1$ and $\hat 1 = \delta_0$; the Fourier transform of a polynomial is a combination of derivatives of $\delta_0$. Compactly supported distributions are tempered, and so are the functions in
$L^p(\RR^n)$ for $1 \leq p \leq \infty$.
\end{example}

\section{Sobolev spaces}\label{distributions:section-sobolev}

\begin{definition}\label{distributions:definition-sobolev}
For $s \in \RR$, $H^s(\RR^n) = \{u \in \mathcal{S}' : \hat u \in L^1_{\mathrm{loc}} \text{ and } \|u\|_s^2 = \int\langle\xi\rangle^{2s}|\hat u(\xi)|^2\,d\xi < \infty\}$, a Hilbert space with the evident inner
product. On the torus $T^n = \RR^n/\ZZ^n$, a distribution $u$ has Fourier coefficients $\hat u(k) = u(e^{-2\pi ik\cdot x})$, $k \in \ZZ^n$, and
$H^s(T^n) = \{u : \|u\|_s^2 = \sum_{k \in \ZZ^n}\langle k \rangle^{2s}|\hat u(k)|^2 < \infty\}$.
\end{definition}

\begin{proposition}\label{distributions:proposition-sobolev-basic}
$H^0 = L^2$, and for a non-negative integer $m$, $H^m$ is the space of $u \in L^2$ whose distributional derivatives of order $\leq m$ lie in $L^2$, with an equivalent norm
$(\sum_{|\alpha| \leq m}\|D^\alpha u\|_2^2)^{1/2}$. The spaces $H^s$ decrease in $s$, $C^\infty$ functions (with compact support, respectively on the torus) are dense in each, and the pairing
$(u, v) \mapsto \int u\bar v$ identifies $H^{-s}$ with the dual of $H^s$.
\end{proposition}

\begin{proof}
$H^0 = L^2$ is Plancherel's theorem (Theorem~\ref{distributions:theorem-fourier-schwartz}), or Parseval's identity on the torus (Example~\ref{functional-analysis:example-fourier}). For integer $m$,
$\mathcal{F}(D^\alpha u) = (2\pi i\xi)^\alpha\hat u$, and $\sum_{|\alpha| \leq m}|(2\pi\xi)^\alpha|^2$ is bounded above and below by multiples of $\langle\xi\rangle^{2m}$. Density follows by truncating and
smoothing the Fourier transform, and duality from the Riesz representation theorem applied to the weighted $L^2$ spaces.
\end{proof}

\begin{theorem}[Sobolev embedding]\label{distributions:theorem-sobolev-embedding}
If $s > k + n/2$, then every $u \in H^s$ (on $\RR^n$ or on $T^n$) has a representative of class $C^k$, and $\|u\|_{C^k} \leq C\|u\|_s$. Hence $\bigcap_sH^s$ consists of the smooth functions.
\end{theorem}

\begin{proof}
On the torus, for $|\alpha| \leq k$,
$\sum_k|(2\pi ik)^\alpha\hat u(k)| \leq \left(\sum_k\langle k\rangle^{2s}|\hat u(k)|^2\right)^{1/2}\left(\sum_k\langle k\rangle^{2(|\alpha| - s)}\right)^{1/2}$ by Cauchy--Schwarz, and the second factor is
finite because $2(s - |\alpha|) > n$ (compare with $\int\langle\xi\rangle^{2(|\alpha|-s)}d\xi$, which converges by polar coordinates). So the Fourier series of $D^\alpha u$ converges absolutely and
uniformly, and $u$ is $C^k$ with the stated bound. On $\RR^n$ the same computation with integrals shows that $\widehat{D^\alpha u} \in L^1$, so $D^\alpha u$ is continuous and bounded by Fourier inversion.
\end{proof}

The embedding theorem trades $n/2$ derivatives in $L^2$ for continuity. With $L^p$ norms one can do better, and the exponents are dictated by scaling: for $u_\lambda(x) = u(\lambda x)$,
$\|u_\lambda\|_q = \lambda^{-n/q}\|u\|_q$ and $\|\nabla u_\lambda\|_p = \lambda^{1 - n/p}\|\nabla u\|_p$, so an inequality $\|u\|_q \leq C\|\nabla u\|_p$ for all $u \in C_c^1(\RR^n)$ can only hold if
$-n/q = 1 - n/p$, that is, $q = p^* = np/(n - p)$, the \emph{Sobolev exponent}, which requires $p < n$.

\begin{theorem}[Gagliardo--Nirenberg--Sobolev]\label{distributions:theorem-gns}
Let $n \geq 2$ and $1 \leq p < n$. For $u \in C_c^1(\RR^n)$, $\|u\|_{n/(n-1)} \leq \|\nabla u\|_1$, and $\|u\|_{p^*} \leq C_{n,p}\|\nabla u\|_p$ with $p^* = np/(n - p)$.
\end{theorem}

\begin{proof}[Sketch of proof]
\emph{Case $p = 1$, $n = 2$.} By the fundamental theorem of calculus in each variable, $|u(x_1, x_2)| \leq \int_\RR|\partial_1u(y_1, x_2)|\,dy_1 =: g_1(x_2)$ and $|u(x_1, x_2)| \leq g_2(x_1)$ similarly. So
$|u(x)|^2 \leq g_1(x_2)g_2(x_1)$, and integrating, $\int|u|^2 \leq \int g_1\int g_2 = \|\partial_1u\|_1\|\partial_2u\|_1 \leq \|\nabla u\|_1^2$. For general $n$ the same bounds give
$|u(x)|^{n/(n-1)} \leq \prod_i g_i^{1/(n-1)}$, where $g_i$ does not depend on $x_i$, and one integrates in $x_1, \ldots, x_n$ in turn, using at each step the generalised H\"older inequality
$\int\prod_{j=1}^{n-1}|h_j| \leq \prod_j\|h_j\|_{n-1}$; this omitted computation gives $\int|u|^{n/(n-1)} \leq \prod_i\|\partial_iu\|_1^{1/(n-1)}$ \cite[Chapter 9]{brezis-functional}.

\emph{From $p = 1$ to $p > 1$.} Apply the case $p = 1$ to $v = |u|^\gamma$ with $\gamma = p(n - 1)/(n - p) > 1$, which is $C^1$ with $|\nabla v| \leq \gamma|u|^{\gamma - 1}|\nabla u|$. By H\"older's inequality
(Proposition~\ref{measure-theory:proposition-holder}) with $p^{\prime} = p/(p - 1)$,
\[
\Bigl(\int|u|^{\gamma n/(n-1)}\Bigr)^{(n-1)/n} \leq \gamma\int|u|^{\gamma - 1}|\nabla u| \leq \gamma\Bigl(\int|u|^{(\gamma - 1)p^{\prime}}\Bigr)^{1/p^{\prime}}\|\nabla u\|_p .
\]
The choice of $\gamma$ makes $\gamma n/(n - 1) = (\gamma - 1)p^{\prime} = p^*$, and dividing gives $\|u\|_{p^*}^{\gamma - p^*/p^{\prime}} \leq \gamma\|\nabla u\|_p$ with $\gamma - p^*/p^{\prime} = 1$.
\end{proof}

For $p = n$ the inequality fails, $p^* = \infty$: in $\RR^2$ the function $\log\log(1/|x|)$, cut off near $|x| = e^{-1}$, has $\nabla u \in L^2$ but is unbounded near $0$. For $p > n$, functions with
$\nabla u \in L^p$ are H\"older continuous (Morrey's inequality \cite[Chapter 9]{brezis-functional}), the $L^p$ analogue of Theorem~\ref{distributions:theorem-sobolev-embedding}.

\begin{theorem}[Rellich]\label{distributions:theorem-rellich}
For $s > t$ the inclusion $H^s(T^n) \to H^t(T^n)$ is a compact operator.
\end{theorem}

\begin{proof}
Let $P_N$ be the truncation $\widehat{P_Nu}(k) = \hat u(k)$ for $|k| \leq N$ and $0$ otherwise, a finite rank operator. For $u \in H^s$,
\[
\|u - P_Nu\|_t^2 = \sum_{|k| > N}\langle k\rangle^{2t}|\hat u(k)|^2 \leq \langle N\rangle^{2(t-s)}\sum_{|k| > N}\langle k\rangle^{2s}|\hat u(k)|^2 \leq \langle N\rangle^{2(t-s)}\|u\|_s^2,
\]
so the inclusion is the norm limit of finite rank operators, hence compact (Proposition~\ref{functional-analysis:proposition-compact-operators}).
\end{proof}

\begin{counterexample}\label{distributions:counterexample-rellich-rn}
Rellich's theorem fails on $\RR^n$: translates $u_j(x) = u(x - j)$ of a fixed nonzero $u \in C_c^\infty$ have constant $H^s$ norms and no subsequence converging in $H^t$, since they are eventually supported
in disjoint sets. Compactness holds for the inclusions $H^s_K \to H^t$ of the subspaces of distributions supported in a fixed compact $K$, which is the form used on manifolds.
\end{counterexample}

\subsection{Poincar\'e inequalities}\label{distributions:subsection-poincare}

A function is not controlled by its derivative alone: adding a constant changes the function and not its derivative. The \emph{Poincar\'e inequalities} say that, once the constants are excluded,
the $L^2$ norm of a function is bounded by the $L^2$ norm of its gradient. There are two ways to exclude the constants: by asking that the function vanish on the boundary, or by asking that its mean
be zero. The inequalities are used constantly in analysis. They make $\int|\nabla u|^2$, the \emph{Dirichlet energy}, a norm, which is the starting point of the variational solution of
the Dirichlet problem (Proposition~\ref{distributions:proposition-dirichlet-weak} below). The best constants are eigenvalues of the Laplacian, a spectral gap that reappears in Hodge theory. We start
with one variable, where the best constants come from Fourier series.

Below, a continuous function $f \colon [a, b] \to \CC$ is \emph{piecewise $C^1$} if there is a subdivision $a = t_0 < t_1 < \cdots < t_m = b$ such that $f$ is $C^1$ on each $[t_{i-1}, t_i]$, with one-sided
derivatives at the ends; $f^{\prime}$ is then defined except at the $t_i$, which does not affect its integrals. A function on $\RR$ is piecewise $C^1$ if it is so on every bounded interval.

\begin{proposition}[Wirtinger's inequality]\label{distributions:proposition-wirtinger}
Let $L > 0$ and let $f \colon \RR \to \CC$ be continuous, piecewise $C^1$ and $L$-periodic, with $\int_0^Lf = 0$. Then
\[
\int_0^L|f|^2 \leq \Bigl(\frac{L}{2\pi}\Bigr)^2\int_0^L|f^{\prime}|^2,
\]
with equality if and only if $f(x) = ae^{2\pi ix/L} + be^{-2\pi ix/L}$ for constants $a, b \in \CC$.
\end{proposition}

\begin{proof}
First let $L = 1$. Write $\hat f(n) = \int_0^1f(t)e^{-2\pi int}\,dt$. Integrating by parts on each interval of $C^1$-ness, the boundary terms cancel, because $fe^{-2\pi int}$ is continuous and
$1$-periodic; so
\[
\widehat{f^{\prime}}(n) = \int_0^1f^{\prime}(t)e^{-2\pi int}\,dt = 2\pi in\int_0^1f(t)e^{-2\pi int}\,dt = 2\pi in\,\hat f(n).
\]
Also $\hat f(0) = \int_0^1f = 0$. By Parseval's identity (Example~\ref{functional-analysis:example-fourier}), applied to $f$ and to the bounded function $f^{\prime}$,
\[
\int_0^1|f^{\prime}|^2 = \sum_{n \in \ZZ}4\pi^2n^2|\hat f(n)|^2 \geq 4\pi^2\sum_{n \neq 0}|\hat f(n)|^2 = 4\pi^2\int_0^1|f|^2,
\]
with equality if and only if $\hat f(n) = 0$ for all $|n| \geq 2$, that is, $f = ae^{2\pi it} + be^{-2\pi it}$ (a function in $L^2$ is determined by its Fourier coefficients). For general $L$ apply
this to $g(t) = f(Lt)$: then $\int_0^L|f|^2 = L\int_0^1|g|^2$, and $g^{\prime}(t) = Lf^{\prime}(Lt)$ gives $\int_0^L|f^{\prime}|^2 = L^{-1}\int_0^1|g^{\prime}|^2$. So
$\int_0^L|f|^2 = L\int_0^1|g|^2 \leq \frac{L}{4\pi^2}\int_0^1|g^{\prime}|^2 = \frac{L^2}{4\pi^2}\int_0^L|f^{\prime}|^2$.
\end{proof}

The extremal functions are the trigonometric functions of the longest period allowed: the inequality says that among periodic functions with mean zero, the slowest oscillation is the cheapest one
in terms of $\int|f^{\prime}|^2$. On an interval the two boundary conditions are handled by extending the function to a periodic one.

\begin{corollary}\label{distributions:corollary-poincare-interval}
Let $f \colon [0, L] \to \CC$ be continuous and piecewise $C^1$.
\begin{enumerate}
\item (\emph{Dirichlet condition.}) If $f(0) = f(L) = 0$, then $\int_0^L|f|^2 \leq (L/\pi)^2\int_0^L|f^{\prime}|^2$, with equality if and only if $f(x) = c\sin(\pi x/L)$.
\item (\emph{Mean zero.}) With $\bar f = \frac1L\int_0^Lf$, $\int_0^L|f - \bar f|^2 \leq (L/\pi)^2\int_0^L|f^{\prime}|^2$, with equality if and only if $f(x) = a + b\cos(\pi x/L)$.
\end{enumerate}
\end{corollary}

\begin{proof}
(1) Let $F$ be the odd extension of $f$ to $[-L, L]$, $F(-x) = -F(x)$, continued $2L$-periodically. Since $f(0) = f(L) = 0$, $F$ is continuous, and it is piecewise $C^1$. Its mean over a
period is $0$ because $F$ is odd, and $\int_{-L}^L|F|^2 = 2\int_0^L|f|^2$, $\int_{-L}^L|F^{\prime}|^2 = 2\int_0^L|f^{\prime}|^2$. Proposition~\ref{distributions:proposition-wirtinger} with period $2L$ gives
the inequality with constant $(2L/2\pi)^2 = (L/\pi)^2$. Equality holds if and only if $F(x) = ae^{\pi ix/L} + be^{-\pi ix/L}$; this is odd exactly when $b = -a$, that is, $F(x) = 2ia\sin(\pi x/L)$.

(2) Let $F$ be the even extension, $F(x) = f(|x|)$ on $[-L, L]$, continued $2L$-periodically; it is continuous because $F(-L) = F(L)$, and piecewise $C^1$. Apply
Proposition~\ref{distributions:proposition-wirtinger} to $F - \bar f$, whose mean over a period is $\frac{1}{2L}\cdot 2\int_0^Lf - \bar f = 0$. The integrals double as in (1), and the extremal functions
$ae^{\pi ix/L} + be^{-\pi ix/L}$ are even exactly when $a = b$, giving $f - \bar f = 2a\cos(\pi x/L)$.
\end{proof}

The constant is the same in both cases, but the extremals differ: for the Dirichlet condition the best function is half a sine wave, which vanishes at both ends; for the mean zero condition it
is half a cosine wave, which is free at the ends. In both cases the constant $(L/\pi)^2$ is the inverse of the smallest positive eigenvalue of $-d^2/dx^2$ on $[0, L]$ with the corresponding
boundary condition: $-(\sin(\pi x/L))^{\prime\prime} = (\pi/L)^2\sin(\pi x/L)$, and similarly for the cosine. In higher dimensions the two versions are the following.

\begin{theorem}[Poincar\'e inequality]\label{distributions:theorem-poincare-dirichlet}
Let $\Omega \subseteq \RR^n$ be open and contained in a slab $\{x : a < x_1 < a + L\}$. Then for every $u \in C^1_c(\Omega)$,
\[
\int_\Omega|u|^2 \leq \Bigl(\frac{L}{\pi}\Bigr)^2\int_\Omega|\partial_1u|^2 \leq \Bigl(\frac{L}{\pi}\Bigr)^2\int_\Omega|\nabla u|^2,
\]
where $|\nabla u|^2 = \sum_j|\partial_ju|^2$.
\end{theorem}

\begin{proof}
Extend $u$ by zero to $\RR^n$; it is $C^1$ with compact support in the open slab. For fixed $x^{\prime} = (x_2, \ldots, x_n)$ the function $f(t) = u(t, x^{\prime})$ is $C^1$ on $[a, a + L]$ and vanishes
near both ends, so Corollary~\ref{distributions:corollary-poincare-interval}(1) (translated to $[a, a + L]$) gives $\int_a^{a+L}|u(t, x^{\prime})|^2\,dt \leq (L/\pi)^2\int_a^{a+L}|\partial_1u(t, x^{\prime})|^2\,dt$.
Integrating over $x^{\prime} \in \RR^{n-1}$ and using Fubini's theorem (Theorem~\ref{measure-theory:theorem-fubini}) gives the first inequality; the second is $|\partial_1u|^2 \leq |\nabla u|^2$.
\end{proof}

Only the width of $\Omega$ in one direction matters: the inequality holds on an infinite strip. Some boundedness is needed, though.

\begin{counterexample}\label{distributions:counterexample-poincare-unbounded}
There is no Poincar\'e inequality $\int|u|^2 \leq C\int|\nabla u|^2$ for $u \in C^1_c(\RR^n)$. Fix $u \neq 0$ and put $u_R(x) = u(x/R)$. Substituting $x = Ry$,
$\int|u_R|^2 = R^n\int|u|^2$ and $\int|\nabla u_R|^2 = R^{n-2}\int|\nabla u|^2$, so the ratio $\int|u_R|^2/\int|\nabla u_R|^2 = R^2\int|u|^2/\int|\nabla u|^2$ is unbounded as $R \to \infty$. Spreading a bump
out over a large region makes it flat without making it small.
\end{counterexample}

\begin{theorem}[Poincar\'e--Wirtinger inequality on a cube]\label{distributions:theorem-poincare-cube}
Let $Q = (0, L)^n$ and $u \in C^1(\bar Q)$, and let $u_Q = L^{-n}\int_Qu$ be its mean. Then
\[
\int_Q|u - u_Q|^2 \leq \Bigl(\frac{L}{\pi}\Bigr)^2\int_Q|\nabla u|^2,
\]
and equality holds for $u(x) = \cos(\pi x_1/L)$.
\end{theorem}

\begin{proof}
We induct on $n$; the case $n = 1$ is Corollary~\ref{distributions:corollary-poincare-interval}(2). Let $n \geq 2$, write $x = (y, t)$ with $y \in Q^{\prime} = (0, L)^{n-1}$ and $t \in (0, L)$, and let
\[
m(t) = L^{1-n}\int_{Q^{\prime}}u(y, t)\,dy
\]
be the mean of $u$ on the slice at height $t$. For each $t$ the function $y \mapsto u(y, t) - m(t)$ has mean zero on $Q^{\prime}$, so it is orthogonal in $L^2(Q^{\prime})$ to the constant $m(t) - u_Q$, and
\[
\int_{Q^{\prime}}|u(y, t) - u_Q|^2\,dy = \int_{Q^{\prime}}|u(y, t) - m(t)|^2\,dy + L^{n-1}|m(t) - u_Q|^2 .
\]
Integrate over $t$. By the induction hypothesis on each slice, the first term contributes at most $(L/\pi)^2\int_Q\sum_{j < n}|\partial_ju|^2$. For the second, $m$ is $C^1$ on $[0, L]$ with
$m^{\prime}(t) = L^{1-n}\int_{Q^{\prime}}\partial_tu(y, t)\,dy$, by differentiation under the integral sign (dominated convergence, Theorem~\ref{measure-theory:theorem-dominated-convergence}, applied to difference
quotients, which are bounded by $\sup|\partial_tu|$ by the mean value theorem). Its mean over $(0, L)$ is $u_Q$, so by Corollary~\ref{distributions:corollary-poincare-interval}(2) and the
Cauchy--Schwarz inequality $|m^{\prime}(t)|^2 \leq L^{1-n}\int_{Q^{\prime}}|\partial_tu(y, t)|^2\,dy$,
\[
L^{n-1}\int_0^L|m(t) - u_Q|^2\,dt \leq L^{n-1}\Bigl(\frac{L}{\pi}\Bigr)^2\int_0^L|m^{\prime}(t)|^2\,dt \leq \Bigl(\frac{L}{\pi}\Bigr)^2\int_Q|\partial_tu|^2 .
\]
Adding the two contributions gives the inequality. For $u = \cos(\pi x_1/L)$ the mean is $0$, $\int_Q|u|^2 = L^n/2$ and $\int_Q|\nabla u|^2 = (\pi/L)^2L^n/2$.
\end{proof}

The argument shows a general principle, the \emph{tensorisation} of Poincar\'e inequalities: on a product, the variance splits into the variance within the slices and the variance of the slice
means, and each is controlled by derivatives in its own directions. So the constant for the cube equals the constant for the interval, independently of the dimension.

\begin{counterexample}\label{distributions:counterexample-poincare-disconnected}
The mean zero version fails on disconnected sets: on the union $\Omega$ of two disjoint open cubes of the same size, the function equal to $1$ on one and $-1$ on the other has mean zero and gradient
zero. On a connected domain made of two cubes joined by a thin tube, the inequality holds, but the best constant becomes large as the tube gets thinner: a function close to $\pm1$ on the two
cubes, changing sign along the tube, has small Dirichlet energy. The best constant of the Poincar\'e--Wirtinger inequality measures how well connected a domain is.
\end{counterexample}

\begin{remark}\label{distributions:remark-poincare-general}
For a bounded connected open set $\Omega$ with Lipschitz boundary, $\int_\Omega|u - u_\Omega|^2 \leq C_\Omega\int_\Omega|\nabla u|^2$ for all $u$ in the Sobolev space of $L^2$ functions with $L^2$
gradient; the proof is by contradiction, using a compactness theorem of Rellich type for such domains \cite[Chapter 9]{brezis-functional}, \cite{evans-gariepy}. The best constants are $1/\lambda$,
where $\lambda$ is the smallest positive eigenvalue of $-\sum_j\partial_j^2$ with the corresponding boundary condition (zero boundary values, respectively zero normal derivative). For convex domains of
diameter $D$, Payne and Weinberger proved the sharp bound $C_\Omega \leq (D/\pi)^2$ \cite{payne-weinberger-1960}, independent of the dimension, and the interval of length $D$ shows that it cannot be
improved. On a compact connected Riemannian manifold the same inequality holds with $1/\lambda_1$, the inverse of the first nonzero eigenvalue of the Laplacian (Chapter~\ref{hodge-theorem}). The
Poincar\'e inequality also has $L^p$ versions, $\|u - u_\Omega\|_p \leq C\|\nabla u\|_p$, and it is related to the Sobolev inequality (Theorem~\ref{distributions:theorem-gns}), which improves the
exponent on the left at the price of a gradient on all of $\RR^n$.
\end{remark}

The Dirichlet version gives the simplest existence theorem for an elliptic equation. Let $\Omega \subseteq \RR^n$ be open. We identify functions on $\Omega$ extended by zero with functions on
$\RR^n$, and let $H^1_0(\Omega)$ be the closure of $C^\infty_c(\Omega)$ in $H^1(\RR^n)$, with the norm $\|u\|_{H^1}^2 = \|u\|_2^2 + \|\nabla u\|_2^2$
(Proposition~\ref{distributions:proposition-sobolev-basic}). Its elements are the $L^2$ functions with $L^2$ distributional derivatives that ``vanish on $\partial\Omega$'' in a weak sense.

\begin{proposition}[Dirichlet principle]\label{distributions:proposition-dirichlet-weak}
Let $\Omega$ be contained in a slab of width $L$, as in Theorem~\ref{distributions:theorem-poincare-dirichlet}. Then:
\begin{enumerate}
\item $\int_\Omega|u|^2 \leq (L/\pi)^2\int_\Omega|\nabla u|^2$ for every $u \in H^1_0(\Omega)$, and $\langle u, v\rangle_\nabla = \int_\Omega\nabla u \cdot \overline{\nabla v}$ is an inner product on $H^1_0(\Omega)$ whose norm
is equivalent to the $H^1$ norm.
\item For every $f \in L^2(\Omega)$ there is a unique $u \in H^1_0(\Omega)$ with $\int_\Omega\nabla u \cdot \overline{\nabla v} = \int_\Omega f\bar v$ for all $v \in H^1_0(\Omega)$. In particular
$-\sum_j\partial_j^2u = f$ as distributions on $\Omega$: $u$ is a \emph{weak solution} of the Dirichlet problem $-\sum_j\partial_j^2u = f$ in $\Omega$, $u = 0$ on $\partial\Omega$.
\end{enumerate}
\end{proposition}

\begin{proof}
(1) Both sides of the inequality are continuous for the $H^1$ norm, and it holds on the dense subspace $C^\infty_c(\Omega)$ by Theorem~\ref{distributions:theorem-poincare-dirichlet}; so it holds on
$H^1_0(\Omega)$. Hence $\|\nabla u\|_2^2 \leq \|u\|_{H^1}^2 \leq (1 + (L/\pi)^2)\|\nabla u\|_2^2$, so $\langle -, - \rangle_\nabla$ is positive definite and its norm is equivalent to the $H^1$ norm. In
particular $H^1_0(\Omega)$, a closed subspace of the Hilbert space $H^1(\RR^n)$, is a Hilbert space for $\langle -, - \rangle_\nabla$.

(2) The functional $\Lambda(v) = \int_\Omega v\bar f$ is linear and bounded on $(H^1_0(\Omega), \langle -, - \rangle_\nabla)$: $|\Lambda(v)| \leq \|f\|_2\|v\|_2 \leq \|f\|_2(L/\pi)\|\nabla v\|_2$ by Cauchy--Schwarz
and (1). By the Riesz representation theorem (Theorem~\ref{functional-analysis:theorem-riesz-representation}) there is a unique $u$ with $\Lambda(v) = \langle v, u\rangle_\nabla$ for all $v$; taking complex
conjugates, this is the stated identity. For $v = \bar\varphi$ with $\varphi \in C^\infty_c(\Omega)$ it reads $\sum_j\int\partial_ju\,\partial_j\varphi = \int f\varphi$, and by the definition of distributional
derivatives (Definition~\ref{distributions:definition-derivative}) the left hand side is $-\sum_j(\partial_j^2u)(\varphi)$.
\end{proof}

The solution is also the unique minimiser of the energy $E(v) = \frac12\int|\nabla v|^2 - \operatorname{Re}\int f\bar v$ on $H^1_0(\Omega)$ (Exercise~\ref{distributions:exercise-dirichlet-energy}). This is
Dirichlet's principle, which Riemann used to construct harmonic and holomorphic functions; Weierstrass objected that a minimiser need not exist, and making the principle rigorous, which Hilbert
began around 1900, was one of the origins of Sobolev spaces. Whether the weak solution is smooth, and takes its boundary values in the classical sense, is a question of regularity, for which the
elliptic estimates of the next section are the basic tool.

\section{Elliptic estimates}\label{distributions:section-elliptic}

\begin{definition}\label{distributions:definition-elliptic-constant}
Let $P(D) = \sum_{|\alpha| \leq m}a_\alpha D^\alpha$ with constant coefficients $a_\alpha \in \CC$. Its \emph{principal symbol} is $\sigma(\xi) = \sum_{|\alpha| = m}a_\alpha(2\pi i\xi)^\alpha$, and $P$ is
\emph{elliptic} if $\sigma(\xi) \neq 0$ for $\xi \neq 0$.
\end{definition}

\begin{example}\label{distributions:example-laplacian}
The Laplacian $\Delta = -\sum_j\partial_j^2$ has symbol $4\pi^2|\xi|^2$ and is elliptic; the heat operator $\partial_t - \Delta$ and the wave operator are not.
\end{example}

\begin{theorem}[Elliptic estimate and regularity]\label{distributions:theorem-elliptic-estimate}
Let $P(D)$ be elliptic of order $m$ with constant coefficients, on the torus $T^n$. For every $s$ there is $C$ with
\[
\|u\|_{s+m} \leq C\left(\|P(D)u\|_s + \|u\|_s\right) \qquad \text{for all } u \in H^{s+m}(T^n).
\]
Moreover, if $u \in H^t(T^n)$ for some $t$ and $P(D)u \in H^s(T^n)$, then $u \in H^{s+m}(T^n)$; in particular $P(D)u \in C^\infty$ implies $u \in C^\infty$.
\end{theorem}

\begin{proof}
On the Fourier side $P(D)$ is multiplication by $p(k) = \sum_{|\alpha| \leq m}a_\alpha(2\pi ik)^\alpha$. By homogeneity and compactness of the unit sphere, $|\sigma(\xi)| \geq c|\xi|^m$ for some $c > 0$, and the
lower-order terms are $O(|\xi|^{m-1})$, so there are $c' > 0$ and $R$ with $|p(k)| \geq c'\langle k\rangle^m$ for $|k| \geq R$. Hence
\[
\sum_{|k| \geq R}\langle k\rangle^{2(s+m)}|\hat u(k)|^2 \leq c'^{-2}\sum_{|k| \geq R}\langle k\rangle^{2s}|p(k)\hat u(k)|^2 \leq c'^{-2}\|P(D)u\|_s^2,
\]
while the finitely many terms with $|k| < R$ are bounded by $\langle R\rangle^{2m}\|u\|_s^2$. This gives the estimate. For regularity, the same computation shows that
$\sum_k\langle k\rangle^{2(s+m)}|\hat u(k)|^2$ is finite whenever $\sum_k\langle k\rangle^{2s}|p(k)\hat u(k)|^2$ is; smoothness follows from
Theorem~\ref{distributions:theorem-sobolev-embedding}, since then $u \in H^s$ for all $s$.
\end{proof}

For operators with variable coefficients the Fourier transform no longer diagonalises $P$, and the estimate is proved by \emph{freezing the coefficients}. Near a point $x_0$, $P$ is close to the
constant-coefficient operator $P_0$ obtained by evaluating its top-order coefficients at $x_0$. The estimate for $P_0$ holds, and the difference $P - P_0$ is either small (top order, near $x_0$) or of
lower order. The errors are absorbed, and the local estimates are patched with a partition of unity. Regularity then follows by applying the estimate to difference quotients. The next subsection
collects the tools: how $\langle\xi\rangle^s$ behaves under translation of $\xi$, multiplication by smooth functions on $H^s$, interpolation between Sobolev norms, and difference quotients.

\subsection{Tools for variable coefficients}\label{distributions:subsection-tools}

\begin{notation}\label{distributions:notation-vector-valued}
Elliptic operators on vector bundles act on vector-valued functions. For $r \geq 1$, $H^s(\RR^n; \CC^r)$ is the space of $u = (u_1, \ldots, u_r)$ with all $u_i \in H^s(\RR^n)$, with
$\|u\|_s^2 = \sum_i\|u_i\|_s^2 = \int\langle\xi\rangle^{2s}|\hat u(\xi)|^2\,d\xi$, where $|\cdot|$ is the Hermitian norm of $\CC^r$. For a matrix $A$, $|A|$ is its operator norm. Products $\varphi u$ of a
matrix-valued function $\varphi$ and a vector-valued $u$ are taken componentwise, $(\varphi u)_i = \sum_j\varphi_{ij}u_j$. We write $\Lambda^su$ for the function with $\widehat{\Lambda^su} = \langle\xi\rangle^s\hat u$, so that
$\|u\|_s = \|\Lambda^su\|_{L^2}$ by Plancherel's theorem (Theorem~\ref{distributions:theorem-fourier-schwartz}(3)); $\Lambda^s$ maps $\mathcal{S}$ to $\mathcal{S}$, since $\langle\xi\rangle^s$ is smooth with derivatives bounded by powers of $\langle\xi\rangle$.
\end{notation}

\begin{lemma}\label{distributions:lemma-peetre}
\begin{enumerate}
\item (\emph{Peetre's inequality}) For $\xi, \eta \in \RR^n$ and $s \in \RR$, $\langle\xi\rangle^s \leq 2^{|s|/2}\langle\xi - \eta\rangle^{|s|}\langle\eta\rangle^s$.
\item (\emph{Convolution bound}) If $k \in L^1(\RR^n)$ and $g \in L^2(\RR^n)$ are nonnegative, then $\|k * g\|_2 \leq \|k\|_1\|g\|_2$.
\item For $\varphi, u \in \mathcal{S}$, $\widehat{\varphi u} = \hat\varphi * \hat u$.
\end{enumerate}
\end{lemma}

\begin{proof}
(1) Since $|\xi|^2 \leq (|\xi - \eta| + |\eta|)^2 \leq 2|\xi - \eta|^2 + 2|\eta|^2$, we get $1 + |\xi|^2 \leq 2(1 + |\xi - \eta|^2)(1 + |\eta|^2)$, that is,
$\langle\xi\rangle^2 \leq 2\langle\xi - \eta\rangle^2\langle\eta\rangle^2$. For $s \geq 0$ raise this to the power $s/2$. For $s < 0$ apply the case $-s > 0$ with $\xi$ and $\eta$ exchanged:
$\langle\eta\rangle^{-s} \leq 2^{-s/2}\langle\eta - \xi\rangle^{-s}\langle\xi\rangle^{-s}$, and rearrange.

(2) By the Cauchy--Schwarz inequality applied to $k(\xi - \eta)^{1/2} \cdot k(\xi - \eta)^{1/2}g(\eta)$,
$(k * g)(\xi)^2 \leq \|k\|_1\int k(\xi - \eta)g(\eta)^2\,d\eta$. Integrating in $\xi$ and using Tonelli's theorem (Theorem~\ref{measure-theory:theorem-fubini}) gives $\|k * g\|_2^2 \leq \|k\|_1^2\|g\|_2^2$.

(3) By Fourier inversion, $\varphi(x) = \int\hat\varphi(\eta)e^{2\pi ix\cdot\eta}\,d\eta$. Substituting and using Fubini's theorem, which applies because $\hat\varphi$ and $u$ are integrable,
$\widehat{\varphi u}(\xi) = \int\hat\varphi(\eta)\int u(x)e^{-2\pi ix\cdot(\xi - \eta)}\,dx\,d\eta = \int\hat\varphi(\eta)\hat u(\xi - \eta)\,d\eta$.
\end{proof}

Peetre's inequality says that the weight $\langle\xi\rangle^s$ changes by a bounded factor when $\xi$ moves by a bounded amount. Multiplying by a smooth function $\varphi$ convolves $\hat u$ with $\hat\varphi$, which
moves frequencies by amounts concentrated near $0$; so it does not change Sobolev regularity.

\begin{proposition}\label{distributions:proposition-multiplication}
Let $\varphi \in \mathcal{S}(\RR^n; \CC^{r \times r})$ and $s \in \RR$, and put $N_t(\varphi) = \int\langle\xi\rangle^{t}|\hat\varphi(\xi)|\,d\xi$ for $t \geq 0$. Then $u \mapsto \varphi u$ extends from $\mathcal{S}$ to a bounded
operator on $H^s(\RR^n; \CC^r)$, and for all $u \in H^s(\RR^n; \CC^r)$:
\begin{enumerate}
\item $\|\varphi u\|_s \leq 2^{|s|/2}N_{|s|}(\varphi)\|u\|_s$;
\item $\|\varphi u\|_s \leq \sup|\varphi|\,\|u\|_s + |s|2^{|s-1|/2}N_{|s-1|+1}(\varphi)\|u\|_{s-1}$.
\end{enumerate}
\end{proposition}

\begin{proof}
First let $u \in \mathcal{S}$. By Lemma~\ref{distributions:lemma-peetre}(3) (componentwise), $\widehat{\varphi u}(\xi) = \int\hat\varphi(\xi - \eta)\hat u(\eta)\,d\eta$.

(1) By Peetre's inequality, $\langle\xi\rangle^s|\widehat{\varphi u}(\xi)| \leq 2^{|s|/2}\int k(\xi - \eta)g(\eta)\,d\eta$ with $k = \langle\cdot\rangle^{|s|}|\hat\varphi|$ and $g = \langle\cdot\rangle^s|\hat u|$. Lemma~\ref{distributions:lemma-peetre}(2)
gives $\|\varphi u\|_s \leq 2^{|s|/2}\|k\|_1\|g\|_2 = 2^{|s|/2}N_{|s|}(\varphi)\|u\|_s$.

(2) Split $\langle\xi\rangle^s\widehat{\varphi u}(\xi) = \int\hat\varphi(\xi - \eta)\langle\eta\rangle^s\hat u(\eta)\,d\eta + \int\hat\varphi(\xi - \eta)\bigl(\langle\xi\rangle^s - \langle\eta\rangle^s\bigr)\hat u(\eta)\,d\eta$. The first term is
$\widehat{\varphi\Lambda^su}(\xi)$, whose $L^2$ norm is $\|\varphi\Lambda^su\|_2 \leq \sup|\varphi|\,\|\Lambda^su\|_2 = \sup|\varphi|\,\|u\|_s$ by Plancherel's theorem. For the second, the gradient of $\zeta \mapsto \langle\zeta\rangle^s$ is
$s\langle\zeta\rangle^{s-2}\zeta$, of length at most $|s|\langle\zeta\rangle^{s-1}$. By the mean value theorem on the segment from $\eta$ to $\xi$, and Peetre's inequality for the points $\zeta$ of this segment,
which satisfy $|\zeta - \eta| \leq |\xi - \eta|$,
\[
\bigl|\langle\xi\rangle^s - \langle\eta\rangle^s\bigr| \leq |s|\,|\xi - \eta|\max_\zeta\langle\zeta\rangle^{s-1} \leq |s|2^{|s-1|/2}\langle\xi - \eta\rangle^{|s-1|+1}\langle\eta\rangle^{s-1} .
\]
Lemma~\ref{distributions:lemma-peetre}(2) bounds the $L^2$ norm of the second term by $|s|2^{|s-1|/2}N_{|s-1|+1}(\varphi)\|u\|_{s-1}$.

\emph{Extension.} $\mathcal{S}$ is dense in $H^s$ (Proposition~\ref{distributions:proposition-sobolev-basic}), so by (1) the map extends uniquely to a bounded operator, and (1), (2) extend by continuity. The
extension is the product of the distribution $u$ with $\varphi$ (Definition~\ref{distributions:definition-derivative}): if $u_k \to u$ in $H^s$ with $u_k \in \mathcal{S}$, then $u_k \to u$ and $\varphi u_k \to \varphi u$ as
distributions, because $|\int v\psi| \leq \|v\|_s\|\bar\psi\|_{-s}$ for test functions $\psi$ (the duality of Proposition~\ref{distributions:proposition-sobolev-basic}).
\end{proof}

Estimate (2) is the one that makes the freezing argument work: up to a term of lower order, multiplication by $\varphi$ costs only the factor $\sup|\varphi|$, which is small when $\varphi$ is small,
however wildly $\varphi$ oscillates.

\begin{example}\label{distributions:example-multiplication}
\begin{enumerate}
\item For $s = 0$ both estimates are the trivial bound $\|\varphi u\|_2 \leq \sup|\varphi|\,\|u\|_2$, up to the constant $N_0(\varphi) \geq \sup|\varphi|$ in (1).
\item Some smoothness of $\varphi$ is needed: multiplication by the indicator function $\mathbf{1}_{[0, \infty)}$ does not preserve $H^1(\RR)$, since it turns a smooth function with $u(0) \neq 0$ into a
function with a jump, which is not in $H^1$ (its derivative contains $u(0)\delta_0$, which is not in $L^2$).
\item The weight in $N_{|s|}(\varphi)$ cannot be replaced by $\sup|\varphi|$ in (1). Multiplication by $e^{2\pi ix\cdot a}$ translates $\hat u$ by $a$. If $\hat u$ is concentrated near $0$, then $\|u\|_s \approx \|u\|_0$, while
$\|e^{2\pi ix\cdot a}u\|_s \approx \langle a\rangle^s\|u\|_0$. So for $s > 0$, multiplication by $\varphi(x) = e^{2\pi ix\cdot a}\psi(x)$, with $\psi \in \mathcal{S}$ equal to $1$ on a large ball, has norm of order $\langle a\rangle^s$ on
$H^s$, although $\sup|\varphi| = \sup|\psi|$ does not depend on $a$. This agrees with Peetre's inequality, and with (2): the oscillation of $\varphi$ only enters the lower-order term.
\end{enumerate}
\end{example}

\begin{lemma}[Interpolation]\label{distributions:lemma-interpolation}
Let $r < t < s$ and $0 < \varepsilon \leq 1$. Then $\|u\|_t \leq \varepsilon\|u\|_s + \varepsilon^{-(t-r)/(s-t)}\|u\|_r$ for $u \in H^s(\RR^n; \CC^r)$. Also $\|\partial^\alpha u\|_t \leq (2\pi)^{|\alpha|}\|u\|_{t+|\alpha|}$.
\end{lemma}

\begin{proof}
Where $\langle\xi\rangle^{s-t} \geq 1/\varepsilon$, $\langle\xi\rangle^t \leq \varepsilon\langle\xi\rangle^s$; elsewhere $\langle\xi\rangle^t = \langle\xi\rangle^{t-r}\langle\xi\rangle^r \leq \varepsilon^{-(t-r)/(s-t)}\langle\xi\rangle^r$. So
$\langle\xi\rangle^t \leq \varepsilon\langle\xi\rangle^s + \varepsilon^{-(t-r)/(s-t)}\langle\xi\rangle^r$ for every $\xi$, and the first claim follows from the triangle inequality in $L^2$ applied to these functions times $|\hat u|$.
The second follows from $|(2\pi i\xi)^\alpha| \leq (2\pi)^{|\alpha|}\langle\xi\rangle^{|\alpha|}$.
\end{proof}

High frequencies are controlled by the strong norm with a small constant, low frequencies by the weak norm with a large one. This is why lower-order terms never matter in elliptic estimates.

\begin{lemma}[Difference quotients]\label{distributions:lemma-difference-quotients}
For $h \in \RR \setminus \{0\}$ and $1 \leq j \leq n$ let $\tau_hu(x) = u(x + he_j)$ and $\Delta_hu = (\tau_hu - u)/h$ (on distributions, $(\tau_hT)(\psi) = T(\tau_{-h}\psi)$).
\begin{enumerate}
\item $\|\tau_hu\|_s = \|u\|_s$ and $\|\Delta_hu\|_s \leq 2\pi\|u\|_{s+1}$.
\item If $u \in H^s$ and there are $A$ and $h_0 > 0$ with $\|\Delta_hu\|_s \leq A$ for all $j$ and all $0 < |h| \leq h_0$, then $u \in H^{s+1}$ and $\|u\|_{s+1}^2 \leq \|u\|_s^2 + nA^2/(4\pi^2)$.
\item For a smooth function $\varphi$, $\varphi\,\Delta_hu = \Delta_h(\varphi u) - (\Delta_h\varphi)\,\tau_hu$; and for $\varphi \in \mathcal{S}$, $N_t(\tau_h\varphi) = N_t(\varphi)$ and $N_t(\Delta_h\varphi) \leq 2\pi N_{t+1}(\varphi)$.
\end{enumerate}
\end{lemma}

\begin{proof}
By the translation rule (Example~\ref{distributions:example-fourier-computations}(3)), $\widehat{\tau_hu}(\xi) = e^{2\pi ih\xi_j}\hat u(\xi)$ and $\widehat{\Delta_hu}(\xi) = \frac{e^{2\pi ih\xi_j} - 1}{h}\hat u(\xi)$.

(1) $|e^{2\pi ih\xi_j}| = 1$, and $|e^{i\theta} - 1| \leq |\theta|$ gives $|e^{2\pi ih\xi_j} - 1|/|h| \leq 2\pi|\xi_j| \leq 2\pi\langle\xi\rangle$.

(2) As $h \to 0$, $|e^{2\pi ih\xi_j} - 1|^2/h^2 \to 4\pi^2\xi_j^2$. By Fatou's lemma (Corollary~\ref{measure-theory:corollary-integral-properties}(2)),
$\int\langle\xi\rangle^{2s}4\pi^2\xi_j^2|\hat u|^2 \leq \liminf_{h \to 0}\|\Delta_hu\|_s^2 \leq A^2$. Summing over $j$ and using $\langle\xi\rangle^{2s+2} = \langle\xi\rangle^{2s}(1 + |\xi|^2)$ gives the claim.

(3) $\Delta_h(\varphi u) = (\tau_h\varphi)\Delta_hu + (\Delta_h\varphi)u$, and $\tau_h\varphi - \varphi = h\Delta_h\varphi$, $u + h\Delta_hu = \tau_hu$; so
$\varphi\,\Delta_hu = \Delta_h(\varphi u) - (\Delta_h\varphi)(u + h\Delta_hu) = \Delta_h(\varphi u) - (\Delta_h\varphi)\tau_hu$. The bounds on $N_t$ follow from the formulas for $\widehat{\tau_h\varphi}$ and $\widehat{\Delta_h\varphi}$ and from (1).
\end{proof}

Part (2) is a converse to (1): a function whose difference quotients are bounded in $H^s$ uniformly in $h$ has one more derivative in $H^s$. In $L^2$ this is the familiar fact that a function with
bounded difference quotients in $L^2$ has a derivative in $L^2$; it fails in $L^\infty$, where bounded difference quotients only give Lipschitz functions, which need not be $C^1$.

\begin{lemma}\label{distributions:lemma-constant-coefficients-rn}
Let $P_0 = \sum_{|\alpha| = m}A_\alpha\partial^\alpha$ with constant matrices $A_\alpha \in \CC^{r \times r}$, $m \geq 1$, and suppose that $\sigma_0(\xi) = \sum_{|\alpha| = m}A_\alpha(2\pi i\xi)^\alpha$ is invertible for $\xi \neq 0$.
Let $c = \min\{|\sigma_0(\xi)v| : |\xi| = 1, |v| = 1\} > 0$. Then for every $s$ and all $u \in H^{s+m}(\RR^n; \CC^r)$,
\[
\|u\|_{s+m} \leq 2^{(m-1)/2}\max(1, c^{-1})\bigl(\|P_0u\|_s + \|u\|_s\bigr).
\]
\end{lemma}

\begin{proof}
The minimum is attained, by compactness, and is positive. By homogeneity $|\sigma_0(\xi)v| \geq c|\xi|^m|v|$. By convexity of $t \mapsto t^m$, $(1 + a)^m \leq 2^{m-1}(1 + a^m)$ for $a \geq 0$, so
$\langle\xi\rangle^{2m}|\hat u(\xi)|^2 \leq 2^{m-1}\bigl(|\hat u(\xi)|^2 + c^{-2}|\sigma_0(\xi)\hat u(\xi)|^2\bigr)$. Since $\widehat{P_0u} = \sigma_0\hat u$ (Definition~\ref{distributions:definition-fourier-tempered}), multiplying by
$\langle\xi\rangle^{2s}$ and integrating gives $\|u\|_{s+m}^2 \leq 2^{m-1}(\|u\|_s^2 + c^{-2}\|P_0u\|_s^2)$, and $\sqrt{a^2 + b^2} \leq a + b$.
\end{proof}

\subsection{Variable coefficients}\label{distributions:subsection-variable-coefficients}

\begin{definition}\label{distributions:definition-elliptic-variable}
Let $U \subseteq \RR^n$ be open. A \emph{linear differential operator} of order $m$ on $U$ is $P = \sum_{|\alpha| \leq m}a_\alpha(x)\partial^\alpha$ with $a_\alpha \in C^\infty(U; \CC^{r \times r})$. Its \emph{principal symbol} is
$\sigma_P(x, \xi) = \sum_{|\alpha| = m}a_\alpha(x)(2\pi i\xi)^\alpha$, and $P$ is \emph{elliptic at $x$} if $\sigma_P(x, \xi)$ is invertible for all $\xi \neq 0$.
\end{definition}

\begin{notation}\label{distributions:notation-pu}
If $u \in H^t(\RR^n; \CC^r)$ has support in a compact set $K \subseteq U$, we put $Pu = \sum_\alpha(\chi a_\alpha)\partial^\alpha u$, where $\chi \in C_c^\infty(U)$ equals $1$ on a neighbourhood of $K$
(Corollary~\ref{real-analysis:corollary-bump-function} and a finite cover of $K$ by balls). Since $\partial^\alpha u$ is supported in $K$, this does not depend on $\chi$, and $Pu \in H^{t-m}$ by
Proposition~\ref{distributions:proposition-multiplication} and Lemma~\ref{distributions:lemma-interpolation}. We say that $u$ is supported in $K$ if its support as a distribution is.
\end{notation}

\begin{theorem}[Local elliptic estimate]\label{distributions:theorem-local-estimate}
Let $P$ be a differential operator of order $m \geq 1$ on $U$, and let $K \subseteq U$ be a compact set at every point of which $P$ is elliptic. For every $s \in \RR$ there is $C$ such that
\[
\|u\|_{s+m} \leq C\bigl(\|Pu\|_s + \|u\|_s\bigr) \qquad \text{for all } u \in H^{s+m}(\RR^n; \CC^r) \text{ supported in } K.
\]
\end{theorem}

\begin{proof}
\emph{Step 1: freezing the coefficients at one point.} Fix $x_0 \in K$, let $P_0 = \sum_{|\alpha| = m}a_\alpha(x_0)\partial^\alpha$ and let $C_0$ be the constant of Lemma~\ref{distributions:lemma-constant-coefficients-rn} for $P_0$. Let
$N$ be the number of multi-indices with $|\alpha| = m$. Choose $\delta > 0$ with $B(x_0, 2\delta) \subseteq U$ and
\[
\eta = \max_{|\alpha| = m}\sup_{B(x_0, 2\delta)}|a_\alpha - a_\alpha(x_0)| \leq \frac{1}{2C_0N(2\pi)^m},
\]
possible by continuity, and $\chi \in C_c^\infty(B(x_0, 2\delta))$ with $0 \leq \chi \leq 1$ and $\chi = 1$ on $B(x_0, \delta)$. Let $v \in H^{s+m}$ be supported in a compact subset of $B(x_0, \delta)$. Then
$Pv - P_0v = \sum_{|\alpha| = m}\varphi_\alpha\partial^\alpha v + \sum_{|\alpha| < m}\chi a_\alpha\partial^\alpha v$ with $\varphi_\alpha = \chi(a_\alpha - a_\alpha(x_0))$, $\sup|\varphi_\alpha| \leq \eta$. By
Proposition~\ref{distributions:proposition-multiplication}(2) and Lemma~\ref{distributions:lemma-interpolation}, $\|\varphi_\alpha\partial^\alpha v\|_s \leq \eta(2\pi)^m\|v\|_{s+m} + C\|v\|_{s+m-1}$, and by
Proposition~\ref{distributions:proposition-multiplication}(1) the lower-order terms are bounded by $C\|v\|_{s+m-1}$. Here and below $C$ denotes constants independent of $v$ (and later of $u$). Hence
\[
\|v\|_{s+m} \leq C_0\bigl(\|P_0v\|_s + \|v\|_s\bigr) \leq C_0\|Pv\|_s + \tfrac12\|v\|_{s+m} + C\|v\|_{s+m-1},
\]
using $\|v\|_s \leq \|v\|_{s+m-1}$. Since $\|v\|_{s+m}$ is finite, $\|v\|_{s+m} \leq 2C_0\|Pv\|_s + C\|v\|_{s+m-1}$.

\emph{Step 2: a partition of unity.} Cover $K$ by finitely many balls $B(x_i, \delta_i/2)$, $x_i \in K$, with $\delta_i$ as in Step 1. Let $\theta_i \in C_c^\infty(B(x_i, \delta_i))$ be nonnegative and equal to $1$ on
$B(x_i, \delta_i/2)$ (Corollary~\ref{real-analysis:corollary-bump-function}), so $\Theta = \sum_i\theta_i \geq 1$ on a neighbourhood $V$ of $K$, and let $\zeta \in C_c^\infty(V)$ be equal to $1$ near $K$. Then
$\psi_i = \zeta\theta_i/\Theta$ is smooth with compact support in $B(x_i, \delta_i)$, and $\sum_i\psi_i = 1$ near $K$. For $u$ supported in $K$, $u = \sum_i\psi_iu$, and $\psi_iu$ is supported in a compact
subset of $B(x_i, \delta_i)$. By the Leibniz rule (which holds for distributions, by Definition~\ref{distributions:definition-derivative}),
$P(\psi_iu) = \psi_iPu + \sum_{|\alpha| \leq m}\sum_{0 < \beta \leq \alpha}\binom{\alpha}{\beta}\chi a_\alpha(\partial^\beta\psi_i)\partial^{\alpha - \beta}u$, and the double sum has $H^s$ norm at most $C\|u\|_{s+m-1}$. With Step 1,
\[
\begin{aligned}
\|u\|_{s+m} &\leq \sum_i\|\psi_iu\|_{s+m} \leq C\sum_i\bigl(\|P(\psi_iu)\|_s + \|\psi_iu\|_{s+m-1}\bigr) \\
&\leq C\bigl(\|Pu\|_s + \|u\|_{s+m-1}\bigr).
\end{aligned}
\]

\emph{Step 3: absorbing the error.} If $m = 1$ this is the claim. If $m \geq 2$, Lemma~\ref{distributions:lemma-interpolation} with $\varepsilon = 1/(2C)$ gives $C\|u\|_{s+m-1} \leq \frac12\|u\|_{s+m} + C^{\prime}\|u\|_s$, and
absorbing $\frac12\|u\|_{s+m}$ gives the estimate.
\end{proof}

\begin{theorem}[Local elliptic regularity]\label{distributions:theorem-local-regularity}
Let $P$ be a differential operator of order $m \geq 1$ on $U$, elliptic at every point of $U$. If $u \in H^t(\RR^n; \CC^r)$ has compact support in $U$ and $Pu \in H^s$, then $u \in H^{s+m}$. In
particular, if $Pu$ is smooth, so is $u$.
\end{theorem}

\begin{proof}
Let $K$ be the support of $u$ and $\rho > 0$ such that $K^{\prime} = \{x : d(x, K) \leq \rho\} \subseteq U$; we compute $P$ with a cut-off equal to $1$ near $K^{\prime}$, and write $Pv = \sum_\alpha b_\alpha\partial^\alpha v$ with
$b_\alpha \in C_c^\infty$ for $v$ supported in $K^{\prime}$. If $t \geq s + m$ there is nothing to prove. Otherwise we show that $u \in H^{t_1}$ with $t_1 = \min(t + 1, s + m)$; repeating, $u \in H^{s+m}$ after finitely many steps.

Let $0 < |h| \leq \rho$ and $1 \leq j \leq n$. Then $\Delta_hu \in H^t \subseteq H^{t_1 - 1}$ is supported in $K^{\prime}$, and Theorem~\ref{distributions:theorem-local-estimate} on $K^{\prime}$ with $s$ replaced by
$t_1 - 1 - m$ gives
\[
\|\Delta_hu\|_{t_1 - 1} \leq C\bigl(\|P\Delta_hu\|_{t_1 - 1 - m} + \|\Delta_hu\|_{t_1 - 1 - m}\bigr).
\]
By Lemma~\ref{distributions:lemma-difference-quotients}(3), applied to $w = \partial^\alpha u$ (difference quotients commute with derivatives), $P\Delta_hu = \Delta_h(Pu) - \sum_\alpha(\Delta_hb_\alpha)\,\partial^\alpha\tau_hu$. We bound
the terms with Lemma~\ref{distributions:lemma-difference-quotients}, Proposition~\ref{distributions:proposition-multiplication}(1) and Lemma~\ref{distributions:lemma-interpolation}, using $t_1 - m \leq s$ and $t_1 - 1 \leq t$:
\[
\begin{aligned}
\|\Delta_h(Pu)\|_{t_1 - 1 - m} &\leq 2\pi\|Pu\|_{t_1 - m} \leq 2\pi\|Pu\|_s, \\
\|(\Delta_hb_\alpha)\,\partial^\alpha\tau_hu\|_{t_1 - 1 - m} &\leq C\,N_{|t_1 - 1 - m| + 1}(b_\alpha)\,\|u\|_{t_1 - 1} \leq C\|u\|_t, \\
\|\Delta_hu\|_{t_1 - 1 - m} &\leq 2\pi\|u\|_{t_1 - m} \leq 2\pi\|u\|_t .
\end{aligned}
\]
The bounds do not depend on $h$ or $j$, so Lemma~\ref{distributions:lemma-difference-quotients}(2) gives $u \in H^{t_1}$. If $Pu$ is smooth, it is in every $H^s$ (it has compact support), so $u \in \bigcap_sH^s$
is smooth by the Sobolev embedding theorem (Theorem~\ref{distributions:theorem-sobolev-embedding}).
\end{proof}

\begin{corollary}\label{distributions:corollary-local-regularity}
Let $P$ be elliptic of order $m \geq 1$ at every point of $U$, and let $u$ be a distribution on $U$ such that $\chi u \in H^t(\RR^n; \CC^r)$ for every $\chi \in C_c^\infty(U)$, for some fixed $t$. If $Pu$ is smooth on $U$,
then $u$ is smooth on $U$.
\end{corollary}

\begin{proof}
We show by induction on $k$ that $\chi u \in H^{t+k}$ for all $\chi \in C_c^\infty(U)$ and all $k \geq 0$. Given $\chi$, choose $\tilde\chi \in C_c^\infty(U)$ equal to $1$ near the support of $\chi$. By the Leibniz rule,
$P(\chi u) = \chi Pu + \sum_{|\alpha| \leq m}\sum_{0 < \beta \leq \alpha}\binom{\alpha}{\beta}a_\alpha(\partial^\beta\chi)\partial^{\alpha - \beta}(\tilde\chi u)$. The first term is smooth with compact support, and by the induction hypothesis
$\tilde\chi u \in H^{t+k}$, so the sum is in $H^{t+k-m+1}$. By Theorem~\ref{distributions:theorem-local-regularity}, $\chi u \in H^{t+k+1}$. Hence $\chi u$ is smooth for every $\chi$, and so is $u$.
\end{proof}

\begin{example}\label{distributions:example-weyl-lemma}
\emph{Weyl's lemma.} If $u \in L^2_{\mathrm{loc}}(U)$ satisfies $\sum_j\partial_j^2u = 0$ as a distribution, then $u$ is (almost everywhere equal to) a smooth harmonic function. Indeed $\chi u \in L^2 = H^0$ for every
$\chi \in C_c^\infty(U)$, and the Laplacian is elliptic, so Corollary~\ref{distributions:corollary-local-regularity} applies. The same argument shows that a locally square integrable weak solution of the Cauchy--Riemann
equation $\frac12(\partial_x + i\partial_y)u = 0$, an operator with principal symbol $\pi i(\xi_1 + i\xi_2) \neq 0$, is a holomorphic function. For the Laplacian the continuity of $u$ is already a strong
statement: a weak solution of $\Delta u = 0$ cannot have the jumps that the wave equation allows (Counterexample~\ref{distributions:counterexample-degenerate} below shows what goes wrong without ellipticity).
\end{example}

\begin{counterexample}\label{distributions:counterexample-degenerate}
Ellipticity at every point is needed, even for operators of order $1$ in one variable. The operator $P = x\frac{d}{dx}$ on $\RR$ has symbol $2\pi ix\xi$, which vanishes at $x = 0$ only. The Heaviside function
$H = \mathbf{1}_{[0, \infty)}$ satisfies $PH = x\delta_0 = 0$, since $(x\delta_0)(\varphi) = \delta_0(x\varphi) = 0$ (Example~\ref{distributions:example-heaviside}). So $PH$ is smooth, while $H$ is not even continuous; the
singularity sits exactly at the point where $P$ fails to be elliptic. Away from $0$, $H$ is smooth, as Corollary~\ref{distributions:corollary-local-regularity} predicts on $U = \RR \setminus \{0\}$.
\end{counterexample}

\begin{remark}\label{distributions:remark-variable-coefficients}
On a compact manifold the local estimates are patched with a partition of unity, and the result is that an elliptic operator $P$ of order $m$ between sections of vector bundles satisfies
$\|u\|_{s+m} \leq C(\|Pu\|_s + \|u\|_s)$, is Fredholm as an operator $H^{s+m} \to H^s$ by Rellich's theorem, and has smooth kernel; this is done in Chapter~\ref{elliptic-operators} and applied to the
Laplacian in Chapter~\ref{hodge-theorem}. A more powerful approach replaces difference quotients by \emph{pseudodifferential operators}, operators of the form
$u \mapsto \int a(x, \xi)\hat u(\xi)e^{2\pi ix\cdot\xi}\,d\xi$ with symbols $a$ more general than polynomials in $\xi$; an elliptic operator then has a \emph{parametrix}, an inverse up to smoothing operators, which
gives estimate and regularity at once \cite[Chapters 7--10]{hormander-pdo1}, \cite[Chapter 6]{warner-foundations}.
\end{remark}

\section{Exercises}\label{distributions:section-exercises}

\begin{exercise}\label{distributions:exercise-delta-limit}
Show that $\varepsilon^{-n}\rho(x/\varepsilon) \to \delta_0$ in $\mathcal{D}'(\RR^n)$ as $\varepsilon \to 0$, for any $\rho \in C_c^\infty$ with $\int\rho = 1$.
\end{exercise}

\begin{solution}
$\int\varepsilon^{-n}\rho(x/\varepsilon)\varphi(x)\,dx = \int\rho(y)\varphi(\varepsilon y)\,dy \to \varphi(0)\int\rho$ by dominated convergence.
\end{solution}

\begin{exercise}\label{distributions:exercise-fundamental-solution}
Show that $E(x) = \frac{1}{2}|x|$ satisfies $-E'' = -\delta_0$ on $\RR$, and that on $\RR^n$ with $n \geq 3$ the function $c_n|x|^{2-n}$ is a fundamental solution of the Laplacian for a suitable constant
$c_n$.
\end{exercise}

\begin{exercise}\label{distributions:exercise-sobolev-sharp}
Show that $H^{n/2}(T^n)$ is not contained in $C^0(T^n)$, so the condition $s > n/2$ in Theorem~\ref{distributions:theorem-sobolev-embedding} is sharp.
\end{exercise}

\begin{exercise}\label{distributions:exercise-heat-kernel}
Using the Fourier transform, solve the heat equation $\partial_tu = -\Delta u$ on $T^n$ with initial datum $u_0 \in H^s$, and show that $u(t) \in C^\infty$ for every $t > 0$.
\end{exercise}

\begin{solution}
$\hat u(k, t) = e^{-4\pi^2|k|^2t}\hat u_0(k)$ solves the equation coefficientwise, and for $t > 0$ the factors decay faster than any power of $\langle k\rangle$, so $u(t) \in H^r$ for all $r$ and
Theorem~\ref{distributions:theorem-sobolev-embedding} applies.
\end{solution}

\begin{exercise}\label{distributions:exercise-poincare-elementary}
Prove a Poincar\'e inequality without Fourier series: if $\Omega$ lies in the slab $\{a < x_1 < a + L\}$ and $u \in C^1_c(\Omega)$, then $\int_\Omega|u|^2 \leq \frac{L^2}{2}\int_\Omega|\partial_1u|^2$.
Compare the constant with the optimal one, $(L/\pi)^2$.
\end{exercise}

\begin{solution}
Extend $u$ by zero. For fixed $x^{\prime}$, $u(x_1, x^{\prime}) = \int_a^{x_1}\partial_1u(t, x^{\prime})\,dt$, so by Cauchy--Schwarz $|u(x_1, x^{\prime})|^2 \leq (x_1 - a)\int_a^{a+L}|\partial_1u(t, x^{\prime})|^2\,dt$.
Integrating over $x_1 \in (a, a + L)$ gives the factor $\int_a^{a+L}(x_1 - a)\,dx_1 = L^2/2$, and then integrate over $x^{\prime}$. Since $\pi^2 \approx 9.87$, the elementary constant $L^2/2$ is about five times
the optimal one.
\end{solution}

\begin{exercise}\label{distributions:exercise-dirichlet-energy}
In the situation of Proposition~\ref{distributions:proposition-dirichlet-weak}(2), show that $E(v) = \frac12\int|\nabla v|^2 - \operatorname{Re}\int f\bar v$ satisfies
$E(v) = \frac12\|v - u\|_\nabla^2 - \frac12\|u\|_\nabla^2$ for $v \in H^1_0(\Omega)$, so that $u$ is the unique minimiser of $E$.
\end{exercise}

\begin{solution}
By the defining property of $u$, $\int f\bar v = \langle u, v\rangle_\nabla$. Hence $E(v) = \frac12\|v\|_\nabla^2 - \operatorname{Re}\langle u, v\rangle_\nabla$, while
$\frac12\|v - u\|_\nabla^2 = \frac12\|v\|_\nabla^2 - \operatorname{Re}\langle u, v\rangle_\nabla + \frac12\|u\|_\nabla^2$. The minimum $-\frac12\|u\|_\nabla^2$ is attained exactly at $v = u$.
\end{solution}

\begin{exercise}\label{distributions:exercise-poincare-scaling}
Let $\Omega \subseteq \RR^n$ be open, and suppose $\int_\Omega|u|^2 \leq C\int_\Omega|\nabla u|^2$ for all $u \in C^1_c(\Omega)$. Show that for $r > 0$ the same holds on $r\Omega$ with constant $r^2C$. Deduce
that the Poincar\'e constant of a ball of radius $r$ is $c_nr^2$ for a constant $c_n$ depending only on $n$.
\end{exercise}

\begin{solution}
For $u \in C^1_c(r\Omega)$ put $w(y) = u(ry)$, in $C^1_c(\Omega)$. Substituting $x = ry$, $\int_{r\Omega}|u|^2 = r^n\int_\Omega|w|^2$ and, since $\nabla w(y) = r(\nabla u)(ry)$,
$\int_{r\Omega}|\nabla u|^2 = r^{n-2}\int_\Omega|\nabla w|^2$. So $\int_{r\Omega}|u|^2 \leq r^nC\int_\Omega|\nabla w|^2 = r^2C\int_{r\Omega}|\nabla u|^2$. Applying this to $r\Omega$ and $r^{-1}$ shows that the best constants
satisfy $C(r\Omega) = r^2C(\Omega)$; take $\Omega$ the unit ball, whose best constant is finite by Theorem~\ref{distributions:theorem-poincare-dirichlet}, and the best constant is invariant under translations.
\end{solution}

\begin{exercise}\label{distributions:exercise-sobolev-algebra}
Let $s > n/2$. Show that $\|uv\|_s \leq C\|u\|_s\|v\|_s$ for $u, v \in \mathcal{S}(\RR^n)$, so that $H^s(\RR^n)$ is an algebra under pointwise multiplication. (Use $\langle\xi\rangle \leq \langle\xi - \eta\rangle + \langle\eta\rangle$.)
\end{exercise}

\begin{solution}
Squaring shows $\langle\xi\rangle \leq \langle\xi - \eta\rangle + \langle\eta\rangle$, as $1 + |\xi|^2 \leq 1 + (|\xi - \eta| + |\eta|)^2 \leq 2 + |\xi - \eta|^2 + |\eta|^2 + 2\langle\xi - \eta\rangle\langle\eta\rangle$. Hence
$\langle\xi\rangle^s \leq 2^s(\langle\xi - \eta\rangle^s + \langle\eta\rangle^s)$. By Lemma~\ref{distributions:lemma-peetre}(3), $\langle\xi\rangle^s|\widehat{uv}(\xi)| \leq 2^s\bigl((a * |\hat v|)(\xi) + (|\hat u| * b)(\xi)\bigr)$ with
$a = \langle\cdot\rangle^s|\hat u|$ and $b = \langle\cdot\rangle^s|\hat v|$. By Lemma~\ref{distributions:lemma-peetre}(2), $\|a * |\hat v|\|_2 \leq \|\hat v\|_1\|a\|_2 = \|\hat v\|_1\|u\|_s$, and by the Cauchy--Schwarz inequality
$\|\hat v\|_1 \leq (\int\langle\xi\rangle^{-2s}d\xi)^{1/2}\|v\|_s$, where the integral is finite because $2s > n$. The other term is symmetric. For $s \leq n/2$ this fails: $H^{n/2}$ contains unbounded functions
(Exercise~\ref{distributions:exercise-sobolev-sharp}), and their powers leave $H^{n/2}$.
\end{solution}

\begin{exercise}\label{distributions:exercise-estimate-not-elliptic}
Show that the local estimate of Theorem~\ref{distributions:theorem-local-estimate} fails for the non-elliptic operator $P = \partial_1$ on $\RR^2$. Find $u_k \in C_c^\infty(\RR^2)$, supported in a fixed compact set, with
$\|Pu_k\|_0 + \|u_k\|_0$ bounded and $\|u_k\|_1 \to \infty$.
\end{exercise}

\begin{solution}
Fix $\varphi \in C_c^\infty(\RR)$, $\varphi \neq 0$, and let $u_k(x, y) = \varphi(x)\varphi(y)e^{2\pi iky}$. Then $\|u_k\|_0 = \|\varphi\|_2^2$ and $\partial_1u_k = \varphi^{\prime}(x)\varphi(y)e^{2\pi iky}$ has the constant norm
$\|\varphi^{\prime}\|_2\|\varphi\|_2$. But $\partial_2u_k = \varphi(x)(\varphi^{\prime}(y) + 2\pi ik\varphi(y))e^{2\pi iky}$ has norm at least $2\pi k\|\varphi\|_2^2 - \|\varphi\|_2\|\varphi^{\prime}\|_2$, and
$\|u_k\|_1 \geq \|\partial_2u_k\|_0/(2\pi)$ by Lemma~\ref{distributions:lemma-interpolation}, which tends to infinity. The functions $u_k$ oscillate in the direction $dy$, where the symbol $2\pi i\xi_1$ of $\partial_1$ vanishes.
\end{solution}
